% conclusion.tex -- Conclusion
\section{Conclusion}

A persistent change in selected traders' local decision rules could
produce a Sell majority without feedback from prices. The fraction
needed depended on the number of interacting traders and the observation
window, while active technical traders allowed financial conditions to
modify later decisions. A Sell majority did not determine the price
decline by itself: trader cash and shares constrained desired orders,
and the dealer's reserve curve determined the prices at which net
trades could complete. Under the corrected constant-product mechanism,
shallow positive reserves accompanied severe declines, while zero
depth permitted only matched trades at a fixed price. Limited resources
therefore did not supply a general protection against a simulated crash.
The financial consequence of widespread selling depended on how those
resources entered execution and price formation.
